\documentclass[12pt]{article}
\usepackage{geometry}
\geometry{margin=1.4in}
\usepackage{setspace}
\setstretch{1.4}
\usepackage{amsmath}

\title{Social Systems as Displacement\\
Institutions, Networks, and the Cost of Collective Deviation}
\author{Diego Rincón}
\date{}

\begin{document}

\maketitle

\begin{abstract}
A society has a ground state $S_0$: the values, norms, and structures it claims to embody. It has an actual state $S^*$: the way it really functions. Displacement $\xi$ is the gap between them. The cost of maintaining this displacement is paid continuously by those who see the gap: cognitive dissonance, moral injury, the labor of adaptation to hypocrisy. The cost of return — fixing the gap, aligning $S^*$ with $S_0$ — is prohibitive: those benefiting from displacement resist, infrastructure is built to lock the deviation in, and the return requires revolutionary change. This is displacement applied to collective systems.
\end{abstract}

\section{Ground States and Actual States}

\subsection{Democracy}

Claimed ground state $S_0$: power flows from the people. Decisions are made by elected representatives accountable to constituents. Laws are applied equally. Voting is free and fair.

Actual state $S^*$: power flows from money. Representatives are accountable to donors more than voters. Laws are applied unequally based on wealth and race. Voting is suppressed where convenient, gerrymandered where needed.

Displacement: $\xi = $ the gap between consent-of-governed and capture-by-capital.

\subsection{Medicine}

Claimed $S_0$: health is the goal. Treatments are chosen based on efficacy. Costs are set to sustain the system, not to extract maximum profit.

Actual $S^*$: profit is the goal. Treatments are chosen based on margin. Costs are set to whatever the market will bear. Life-saving drugs are priced beyond reach.

Displacement: $\xi = $ profit motive vs. health motive.

\subsection{Education}

Claimed $S_0$: knowledge and critical thinking are developed. Students are prepared to understand the world and contribute to society.

Actual $S^*$: credentialism and debt are the true products. Students are trained to obey, not think. The institution's survival matters more than the student's learning.

Displacement: $\xi = $ education vs. credential mill.

\subsection{Law}

Claimed $S_0$: law is blind. Justice is impartial. The innocent are protected, the guilty are punished fairly.

Actual $S^*$: law serves power. Justice depends on wealth. The poor are punished harshly, the rich go free.

Displacement: $\xi = $ impartial justice vs. justice as power.

\section{Who Pays the Cost of Displacement?}

The person who sees the gap pays the cost $D(\xi)$.

A doctor who knows the drug she prescribes is chosen for profit, not efficacy, pays cognitive dissonance and moral injury every day. The cost $D(\xi)$ includes:
- Mental effort to reconcile the gap
- Suppression of her own judgment to follow protocol
- Grief and guilt over harmed patients
- Risk of speaking out (professional retaliation)

A teacher who knows students are being credentialed, not educated, pays the same cost. So does a lawyer, a politician, a parent in a system they see is broken.

The cost $D(\xi)$ is paid by the conscious. The unconscious (those who don't see the gap) pay nothing. This creates a perverse incentive: be blind, and you're comfortable. See clearly, and you suffer.

\section{The Maintenance Cost}

To keep displacement stable — to prevent the gap from becoming visible — institutions invest in:

\textbf{Mythology:} Constant reinforcement of the claimed ground state. Propaganda, marketing, public relations. "Our mission is to serve the people." "We're committed to justice." These narratives must be maintained even as they're violated.

\textbf{Selection:} Promote people who don't see the gap (the careless, the ambitious, the morally flexible). Marginalize those who see it clearly (the whistleblowers, the ethicists, the prophets).

\textbf{Bureaucracy:} Create enough process and layers that no one person can see the full displacement. Distribute responsibility so that everyone can claim they were just following orders.

\textbf{Normalization:} Make the displacement so routine that it's invisible. "That's just how healthcare works." "All politicians do this." "The system is broken, but what can you do?"

\textbf{Punishment:} Silence those who speak. Professional consequences for doctors who say "our system harms patients." Legal threats for lawyers who reveal bias. Defunding for schools that teach actual thinking.

The total maintenance cost $C_{\text{maint}}(\xi, \tau)$ is the investment in keeping the gap invisible and stable.

\section{The Return Cost}

To return $S^*$ to $S_0$ — to align actual practice with claimed values — is catastrophically expensive.

Suppose a healthcare system claims to serve health but actually serves profit. To return to health:
1. Restructure pricing (losses to hospitals and pharma, gains to patients)
2. Retrain staff (doctors must relearn to choose based on efficacy, not margin)
3. Rebuild supply chains (efficient generics replace lucrative branded drugs)
4. Redistribute power (patients gain voice, corporations lose control)
5. Face the reality of past harms (acknowledge all the people harmed for profit)

Each step is resisted by those profiting from displacement. The return cost includes:
- Lost income for those benefiting from the current system
- Restructuring overhead (retraining, rebuilding, redesign)
- Psychological cost (acknowledgment of harm, guilt, loss of identity as "the good guys")
- Uncertainty (the new system is untested; outcomes are unknown)

The return cost is super-linear in the time displaced. A system displaced for 50 years is harder to return than one displaced for 5 years. Infrastructure is locked in. Generations have built careers on the deviation. Reversing it requires extraordinary effort.

\section{Why Systems Stay Displaced}

Because return is so expensive, systems remain displaced for centuries.

Slavery in America: displaced from the ground state of human dignity for 400 years. The return — emancipation, reconstruction, civil rights — has cost lives, wars, ongoing conflict. And it's still not complete. The return cost continues to be paid.

Colonialism: a civilization claimed to bring enlightenment but actually extracted resources and destroyed cultures. The return cost — decolonization, reparations, restoration — has been resisted for 70+ years and is still incomplete.

Gender inequality: claimed ground state is equality; actual state is systematic subordination. Return cost is so high (restructure all of labor, rewrite legal systems, psychologically detoxify men from hierarchy) that progress is glacial.

These systems persist not because no one sees the displacement. They persist because the return cost is so high that even those who see it clearly often calculate it's cheaper to adapt and survive within the displacement than to pay for return.

\section{The Three Levers}

If a society wants to minimize net cost $g(\tau) = C_{\text{maint}} \cdot \tau + C_{\text{return}}(\tau)$, it can pull three levers:

\subsection{1. Reduce the Displacement $\xi$}

Keep the actual state close to the claimed state. Don't make promises you don't intend to keep. Align institutions with their stated values.

This requires:
- Transparency (make the gap visible)
- Accountability (consequences for hypocrisy)
- Continuous alignment (regular checks and course corrections)

Most institutions actively resist this. Visibility is costly. Accountability threatens power.

\subsection{2. Lower the Maintenance Cost}

Rather than investing in mythology and propaganda, invest in actual change. It costs less in the long term to fix the system than to keep lying about it.

This is counterintuitive but true: the resources spent on maintaining a false narrative (PR, internal propaganda, silencing dissidents, selective hiring) often exceed the cost of actually reducing the gap.

A healthcare system that moved toward actual health optimization would save money compared to the current system of marketing, legal defense, and lobbying to protect profitable harms.

\subsection{3. Distribute the Return Cost}

Rather than a sudden revolutionary change (where return cost is catastrophic), make incremental changes. Spread the cost over time so no single group bears the full burden.

This requires:
- Honest acknowledgment of displacement
- Graduated timeline for return
- Transitional support for those harmed by current system and those who must change
- Shared sacrifice (those benefiting from displacement accept losses; those harmed by it accept slower change)

\section{Why Return Fails}

Most attempts at social return fail because:

\subsection{Incomplete Return}

The system moves partway back toward $S_0$ but stops short. Slavery was partially returned (emancipation) but not fully (structural racism persists). Gender inequality was partially returned (voting rights, legal equality) but not fully (economic gaps, violence, reproductive autonomy remain contested).

The system settles at a new, slightly closer, but still displaced state. Comfortable for those in power, still costly for those who see the gap.

\subsection{Reversing Return (Backlash)}

Once a return begins, those who benefited from displacement resist violently. They feel loss. They re-tighten the mythology. They select more aggressively for the unconscious.

The return cost becomes unsustainable. The society gives up and slides back toward the original displacement.

We see this now: civil rights gains reversed, environmental protections dismantled, reproductive rights stripped. Return attempted; backlash succeeded.

\subsection{The Return Trap}

Once a system is aware of its displacement, and returns are attempted, it faces a choice:
- Complete return (acknowledge all harms, restructure everything, pay the full cost)
- Partial return (compromise, acknowledge some harms, reach a new stable displacement)

Partial return is politically easier. But it leaves the gap alive. And those who see it clearly are still paying the cost $D(\xi)$.

The result: ongoing low-grade conflict, moral injury, and eventual another return attempt (which backlashes and fails).

\section{Sustainable Displacement}

The only stable social systems are those where:
1. The displacement is small ($\xi$ is minimal)
2. Most people don't see the gap (the unconscious majority lives comfortably)
3. Those who do see it are few enough to suppress or burn out

Or:

1. The displacement is acknowledged
2. It is framed as necessary (we can't be perfect, but this gap is worth the cost)
3. The cost is shared fairly
4. Return is continuous, not revolutionary

The second path is rarer. Most societies choose the first: unconsciousness through mythology, propaganda, and punishment of clarity.

\section{The Path}

For a society that wants to minimize suffering, the mathematics are clear:

$g(\tau)$ grows with every year of displacement. The system must choose: pay the return cost now (when it's lower), or pay it later (when it's higher) or never (and pay an infinite maintenance cost).

Pay the cost now. Acknowledge the gap. Return toward the ground state. Distribute the burden fairly. Make it continuous policy, not revolutionary upheaval.

Those profiting from displacement will resist. But the mathematics don't care about their preference. Every year of delay makes return more expensive and the social cost higher.

The only sustainable path is alignment. Keep the gap small. Make return continuous policy. Pay the cost in increments rather than in sudden revolutions.

This requires the leadership to choose the long-term survival of the system over short-term comfort for the powerful.

It is rare. But it is the only path to stability.

\end{document}
